\section{Causal Tracing}
\label{apd:causal-tracing}

\subsection{Experimental Settings}
\label{apd:causal-tracing-details}

Note that, in by-layer experimental results, layers  are numbered from 0 to $L-1$ rather than $1$ to $L$.

In \reffig{avg-1000-trace} and \reffig{trace-mlp-disabled} we evaluate mean causal traces over a set of 1000 factual prompts that are known by GPT-2 XL, collected as follows.  We perform greedy generation using facts and fact templates from \cfd, and we identify predicted text that names the correct object $o^c$ before naming any other capitalized word.  We use the text up to but not including the object $o^c$ as the prompt, and we randomly sample 1000 of these texts.  In this sample of known facts, the predicted probability of the correct object token calculated by GPT-2 XL averages 27.0\%.

In the corrupted run, we corrupt the embeddings of the token naming the subject $s$ by adding Gaussian noise $\epsilon \sim \mathcal{N}(0; \nu)$, where $\nu = 3\sigma_t$ is set to be three times larger than the observed standard deviation $\sigma_t$ of token embeddings as sampled over a body of text.
For each run of text, the process is repeated ten times with different samples of corruption noise.  On average, this reduces the correct object token score to 8.47\%, less than one third the original score.

When we restore hidden states from the original run, we substitute the originally calculated values from the same layer and the same token, and then we allow subsequent calculations to proceed without further intervention.  For the experiments in \reffig{infoflow} (and the purple traces throughout the appendix), a single activation vector is restored.  Naturally, restoring the last vector on the last token will fully restore the original predicted scores, but our plotted results show that there are also earlier activation vectors at a second location that also have a strong causal effect: the average maximum score seen by restoring the most impactful activation vector at the last token of the subject is 19.5\%.  In \reffig{infoflow}j where effects are bucketed by layer, the maximum effect is seen around the 15th layer of the last subject token, where the score is raised on average to 15.0\%.


\subsection{Separating MLP and Attn Effects}
\label{apd:mlp-and-attn-trace}

When decomposing the effects into MLP and Attn lookups, we found that restoring single activation vectors from individual MLP and individual Attn lookups had generally negligible effects, suggesting the decisive information is accumulated across layers. Therefore for MLP and Attn lookups, we restored runs of ten values of $\smash{\atl{m}{l}_i}$ (and $\smash{\atl{a}{l}_i}$, respectively) for an interval of layers ranging from $[l_* - 4, ..., l_* + 5]$ (clipping at the edges), where the results are plotted at layer $l_*$.  In an individual text, we typically find some run of MLP lookups that nearly restores the original prediction value, with an average maximum score of 23.6\%.  \reffig{avg-1000-trace}b buckets averages for each token-location pair, and finds the maximum effect at an interval at the last entity token, centered at the the 17th layer, which restores scores to an average of 15.0\%.   For Attn lookups (\reffig{avg-1000-trace}c), the average maximum score over any location is 19.4\%, and when bucketed by location, the maximum effect is centered at the 32nd layer at the last word before prediction,  which restores scores to an average of 16.5\%.

\reffig{apd-ct-line-plots} shows mean causal traces as line plots with 95\% confidence intervals, instead of heatmaps.


\begin{figure*}[h]
\centering
\includegraphics[width=\textwidth]{figs/appendix-lineplot-causal-trace-crop.pdf}%
\caption{\small Mean causal traces of GPT-XL over a sample of 1000 factual statements, shown as a line plot with 95\% confidence intervals.  (a) Shows the same data as \reffig{infoflow}j as a line plot instead of a heatmap; (b) matches \reffig{infoflow}k; (c) matches \reffig{infoflow}m. The confidence intervals confirm that the distinctions between peak and non-peak causal effects at both early and late sites are significant.}
\lblfig{apd-ct-line-plots}
\end{figure*}

\subsection{Traces of EleutherAI GPT-NeoX (20B) and GPT-J (6B) and smaller models}

\begin{figure}[!ht]
\centering
\includegraphics[width=1\textwidth]{figs/appendix-gptj-neo-causal-trace-crop.pdf}
\vspace{-10pt}%
\caption{\small (a, b, c) Causal traces for GPT-NeoX (20B) and (d, e, f) Causal traces for GPT-J (6B).}
\lblfig{apd-causaltrace-gptj-neo}
\end{figure}
\begin{figure*}
\centering
\includegraphics[width=\textwidth]{figs/appendix-lineplot-trace-sizes-crop.pdf}%
\caption{\small Comparing mean causal traces across a wide range of different model sizes.  (Compare to \reffig{apd-ct-line-plots}.)  GPT-medium (a, b, c) has 334 million parameters, GPT-large (d, e, f) has 774 million parameters, and NeoX-20B (g, h, i) has  20 billion parameters.  In addition, NeoX has some architectural variations. Despite the wide range of differences, a similar pattern of localized causal effects is seen across models.  Interestingly, for very large models, some effects are stronger. For example, hidden states before the last subject token have negative causal effects instead of merely low effects, while hidden states at early layers at the last subject token continue to have large positive effects, continuing to implicate MLP.  Also, attention modules with strong causal effects appear earlier in the stack of layers.}
\lblfig{apd-ct-line-plot-sizes}
\end{figure*}

We conduct the causal trace experiment using on GPT-NeoX (20 billion parameters) as well as GPT-J (6 billion parameters).  For GPT-NeoX we adjust the injected noise to $\nu = 0.03$ and in GPT-J we use $\nu = 0.025$ to match embedding magnitudes.  We use the same factual prompts as GPT-2 XL, eliminating cases where the larger models would have predicted a different word for the object. Results are shown in \reffig{apd-causaltrace-gptj-neo}.
GPT-NeoX and GPT-J differ from GPT-2 because they have has fewer layers (44 and 28 layers instead of 48), and a slightly different residual structure across layers.  Nevertheless, the causal traces look similar, with an early site with causal states concentrated at the last token of the subject, a large role for MLP states at that site.  Again, attention dominates at the last token before prediction.

There are some differences compared to GPT-2.   The importance of attention at the first layers of the last subject token is more apparent in GPT-Neo and GPT-J compared to GPT-2,   suggesting that the attention parameters may be playing a larger role in storing factual associations.  This concentration of attention at the beginning may also be due to fewer layers in the Eleuther models: attending to the subject name must be done in a concentrated way at just a layer or two, because there are not enough layers to spread out that computation in the shallower model.
The similarity between the GPT-NeoX and GPT-J and GPT-2 XL traces helps us to understand why ROME continues to work well with higher-parameter models, as seen on our experiments in altering parameters of GPT-J.

To examine effects over a wide range of scales, we also compare causal traces for smaller models GPT-2 Medium and GPT-2 Large.  These smaller models are compared to NeoX-20B in \reffig{apd-ct-line-plot-sizes}.  We find that across sizes and architectural variations, early-site MLP modules continue to have high indirect causal effects at the last subject token, although the layers where effects peak are different from one model to another.

\subsection{Causal Tracing Examples and Further Insights}
\label{apd:causal-tracing-variations}

We include further examples of phenomena that can be observed in causal traces. \reffig{apd-causaltrace-1} shows typical examples across different facts. \reffig{apd-causaltrace-2} discusses examples where decisive hidden states are not at the \textit{last} subject token.  \reffig{apd-trace-mlp-detail} examines traces at an individual token in more detail. 

We note that causal tracing depends on a corruption rule to create baseline input for a model that does not contain all the information needed to make a prediction.  Therefore we ask: are Causal Tracing results fragile if the exact form of the corruption changes?  We test this by expanding the corruption rule: even when additional tokens after the subject name are also corrupted, we find that the results are substantially the same. \reffig{apd-causaltrace-xt} shows causal traces with the expanded corruption rule.
\reffig{apd-ct-line-plots-xt} similarly shows line plots with the expanded corruption rule.

We do find that the noise must be large enough to create large average total effects.  For example, if noise with variance that is much smaller is used (for example if we set $\sigma=\sigma_t$), average total effects become very small, and the small gap in the behavior between clean runs and corrupted run makes it difficult discern indirect effects of mediators. Similarly, if we use a uniform distribution where components range in $\pm 3\sigma$, effects large enough for causl tracing but smaller than a Gaussian distribution.

If instead of using spherical Gaussian noise, we draw noise from $\mathcal{N}(mu, \Sigma)$ where we set $\mu = \mu_t$ and $\Sigma_ = \Sigma_t$ to match the observed distribution over token embeddings, average total effects are also strong enough to perform causal traces.  This is shown in \reffig{apd-ct-line-plot-noises}.

Furthermore, we investigate whether Integrated Gradients (IG)~\citep{sundararajan2017axiomatic} provides the same insights as Causal Tracing.  We find that IG is very sensitive to local features but does not yield the same insights about large-scale global logic that we have been able to obtain using causal traces. \reffig{apd-ig} compares causal traces to IG saliency maps.


\renewcommand{\floatpagefraction}{.99}%
\begin{figure}[!ht]
\centering
\includegraphics[width=\textwidth]{figs/appendix-causaltrace-1-crop.pdf}%
\caption{\small Further examples of causal traces showing appearance of the common lookup pattern on a variety of different types of facts about people and other kinds of entities.  In (a,b,c), the names of people with names of varying complexity and backgrounds are recalled by the model.  In each case, the MLP lookups on the last token of the name are decisive.  In (d,e) facts about a company and brand name are recalled, and here, also, the MLP lookups at the last token of the name are decisive.}
\lblfig{apd-causaltrace-1}
\end{figure}%
\renewcommand{\floatpagefraction}{.8}
\begin{figure*}
\centering
\includegraphics[width=\textwidth]{figs/appendix-causaltrace-2-crop.pdf}
\caption{\small Causal traces show that the last token of the subject name is not always decisive.  (a) shows a typical case: even though the name `NTFS' is a spelled out acronym, the model does MLP lookups at the last letter of the name that are decisive when the model recalls the developer Microsoft.  However, in a very similar sentence (b), we can see that the last words of `Windows Media Player' are \emph{not} decisive; the first word `Windows' is the token that triggers the decisive lookup for information about the manufacturer.  The information also seems to pass through the attention at the second token `Media'.  Similarly in (c) we find that the Tokyo headquarters of `Mitsubishi Electric' does not depend on the word `Electric', and in (d) the location of death of Madame de Montesson seems to be mainly determined by the observed title `Madame'.  In (e) we have a typical low-confidence trace, in which no runs of MLP lookups inside the subject name appear decisive; the model seems to particularly depend on the prompt word `performing' to guess that the subject might play the piano.}
\lblfig{apd-causaltrace-2}
\end{figure*}
\begin{figure*}
\centering
\includegraphics[width=\textwidth]{figs/appendix-ct-xt-crop.pdf}
\caption{\small Causal traces in the presence of additional corruption.  Similar to \reffig{apd-causaltrace-1}, but instead of corrupting only the subject token, these traces also corrupt the token after the subject.  Causal effects are somewhat reduced due to the the model losing some ability to read the relation between the subject and object, but these traces continue to show concentrated causal effects at the last token of the subject even when the last token is not the last token corrupted.  Causal effects of MLP layers at the last subject token continues to be pronounced.}
\lblfig{apd-causaltrace-xt}
\end{figure*}
\begin{figure*}
\includegraphics[width=\textwidth]{figs/appendix-lineplot-noise-variant-crop.pdf}%
\caption{\small Comparing different noise choices. (Compare to \reffig{apd-ct-line-plots}, where noise is chosen as a $3\sigma_t$ spherical Gaussian, where $\sigma_t$ is measured to match the observed spherical variance over tokens.)  In a, b, c we we draw noise from a multivariate Gaussian $\mathcal{N}(\mu; \Sigma)$ where $\mu$ and $\Sigma$ are chosen to match the observed mean and covariance over a sample of tokens.  In d, e, f we draw noise from a uniform distribution in the range $\pm 3\sigma$ instead of a Gaussian distribution.  In both cases, the average total effects measured between the clean run and the corrupted run are large enough to measure causal traces, but the effects are smaller than the choice of $3\sigma_t$ used in the main paper.}
\lblfig{apd-ct-line-plot-noises}
\end{figure*}
\begin{figure*}
\centering
\includegraphics[width=\textwidth]{figs/appendix-trace-mlp-detail-crop.pdf}
\caption{\small Detail view of causal traces, breaking out a representative set of individual cases from the 1000 factual statements that are averaged in \reffig{trace-mlp-disabled}.  Shows the causal trace at a specific subject token, with and without MLP disabled, as described in \refsec{causal-tracing}.  In every case, the token tested is highlighted in a red box.  In (a,b,c,d,e) cases are shown that fit the typical pattern: Restoring individual hidden states at a range of layers has a strong decisive average causal effect at the last token of the subject.  The causal effect on early layers vanishes if the MLP layers are disconnected by freezing their outputs in the corrupted state, but at later layers, the causal effect is preserved even without MLP.   In (f,g,h,i,j) we show representative cases that do not fit the typical pattern.  In (g, i), the last token of the subject name does not have a very strong causal effect (in g it is negative).  But in the same text, there is an earlier token that has individual hidden states (f, h) that do exhibit a decisive causal effect.  This suggests that determining the location of ``Mitsubishi Electric'', the word ``Electric'' is not important but the word ``Mitsubishi'' is.  Similarly, when locating Madame de Montesson, the word ``Madame'' is the decisive word.  (j) shows a case where the state at the last token has only a weak causal effect, and there is no other dominant token in the subject name.
}
\lblfig{apd-trace-mlp-detail}
\end{figure*}
\begin{figure*}
\centering
\includegraphics[width=\textwidth]{figs/lineplot-causaltrace-xt-crop.pdf}%
\caption{\small Similar to \reffig{apd-ct-line-plots}, but with an additional token corrupted after the subject token, as in \reffig{apd-causaltrace-xt}.  We observe that the emergence of strong early-site causal effects at the MLP modules is systematic and appears under a different corruption scheme, confirming that importance of the last subject token is apparent even when the last subject token is never the last token corrupted.}
\lblfig{apd-ct-line-plots-xt}
\end{figure*}
\begin{figure*}
\centering
\includegraphics[width=\textwidth]{figs/appendix-ig-crop.pdf}
\caption{\small Integrated gradients saliency maps, visualizing the same cases as in \reffig{apd-causaltrace-1}.  Here we compare Causal Tracing to the method of Integrated Gradients~\citep{sundararajan2017axiomatic}.  Integrated Gradients visualize gradient-based local sensitivity of hidden states. Here we compute IG using 50 steps of Gauss-Legendre quadrature on gradients of individual hidden states $\smash{h_t^{(l)}}$, or $\smash{m_t^{(l)}}$ (for MLP), or $\smash{a_t^{(l)}}$ (for Attn), with respect to the predicted output token; we plot the norm of the integrated gradient at each state.  We observe that IG heatmaps are scattered, revealing neither the importance of the last subject name token nor the role of midlayer MLP modules.}
\lblfig{apd-ig}
\end{figure*}